\subsection{Validação e Aplicação}

Como parte da validação preliminar, foi realizada uma análise manual de postagens com base em checagens da agência Aos Fatos. O corpus de referência é populado a partir do arquivo \texttt{Analise manual Aos fatos\_v2.xlsx} e armazenado na coleção ChromaDB \texttt{checks\_v2\_rewrite\_cosine\_sentence\_transformer}. As tabelas abaixo apresentam os resultados agregados por parlamentar e partido.

\begin{table}[h!]
    \centering
    \caption{Resultados Preliminares do IFS por Parlamentar}
    \begin{tabular}{rrrr}
\toprule
Ano & IFS médio & Total de publicações & Quantidade de parlamentares \\
\midrule
2023 & 0.7 & 12 & 8 \\
2024 & 0.7 & 19 & 13 \\
\bottomrule
\end{tabular}

\end{table}

\begin{table}[h!]
    \centering
    \caption{Resultados Preliminares do IFS por Partido}
    \begin{tabular}{lrrr}
\toprule
Partido & Quantidade de parlamentares & Total de publicações & \% de parlamentares \\
\midrule
PL & 9 & 17 & 50.0 \\
PSOL & 2 & 3 & 11.1 \\
Avante  & 1 & 2 & 5.6 \\
Novo & 1 & 1 & 5.6 \\
MDB & 1 & 1 & 5.6 \\
PSB  & 1 & 1 & 5.6 \\
Podemos & 1 & 4 & 5.6 \\
Repúblicanos & 1 & 1 & 5.6 \\
União Brasil & 1 & 1 & 5.6 \\
\bottomrule
\end{tabular}

\end{table}

\subsection{Protocolo Experimental}

Embora métricas agregadas ainda não tenham sido reportadas, recomendamos o seguinte protocolo:

\begin{itemize}
    \item \textbf{Divisão dos dados:} 70/15/15 treino/validação/teste conforme definido em \texttt{config.py}.
    \item \textbf{Métricas:} Precisão, recall, F1-score, ROC-AUC e curvas de calibração (diagramas de confiabilidade) devido às saídas probabilísticas.
    \item \textbf{Ablações:} Comparar a cabeça MLP com uma cabeça linear; avaliar configurações com recuperação habilitada versus desabilitada; estudar perda focal contra entropia cruzada padrão.
    \item \textbf{Varreduras de hiperparâmetros:} Explorar taxas de dropout (0,3--0,6), larguras de camadas ocultas e parâmetros de perda focal.
\end{itemize}

\subsection{Análise da Distribuição}

A análise da distribuição do IFS, incluindo histogramas e boxplots, será detalhada nesta seção para ilustrar a dispersão e a concentração dos valores do índice entre os parlamentares.

\subsection{Notas de Implementação}

\begin{itemize}
    \item \textbf{Gerenciamento de dispositivo:} \texttt{config.DEVICE} escolhe automaticamente CUDA quando disponível.
    \item \textbf{Serialização:} Checkpoints podem armazenar um \texttt{state\_dict} bruto ou um dicionário contendo \texttt{model\_state\_dict}; ambos os formatos são suportados.
    \item \textbf{Reprodutibilidade:} Sementes são definidas via \texttt{config.RANDOM\_SEED}, mas determinismo de ponta a ponta requer semear NumPy, PyTorch CUDA e data loaders.
    \item \textbf{Dependências:} Pacotes-chave listados em \texttt{requirements.txt} incluem \texttt{transformers}, \texttt{torch}, \texttt{chromadb}, \texttt{pandas} e \texttt{sentence-transformers}.
\end{itemize}

\subsection{Limitações}

\begin{itemize}
    \item \textbf{Cobertura de dados:} O desempenho depende da cobertura de alegações verificadas e postagens anotadas; deriva de domínio pode prejudicar a precisão.
    \item \textbf{Latência:} Recuperação e inferência do transformador incorrem em sobrecarga computacional; processamento em lote e aceleração são essenciais para uso em tempo real.
    \item \textbf{Lacunas de avaliação:} A falta de benchmarks publicados limita a comparação com trabalhos anteriores.
\end{itemize}
